Everything is computed in Gauge Theory Gravity (Definition \claimref{obj.synthesis.gtg-flat-spacetime-gauge-fields}) and in the conformal model $\Cl(2,4)$: tractors are vectors $\sigma\nInf + \mu + \rho\,\nO$ with $\nO\cdot\nInf = 1$, the tractor connection is the gauge field $\Omega_{T}(a) = \Omega(a) - \nInf\wedge\hb^{-1}(a) + \nO\wedge\Sch(\hb^{-1}(a))$ of Theorem \claimref{synthesis.gtg-tractor-gauge-field-reproduces-tractor-connection}, acting by the commutator, $D_{a}V = a\cdot\nabla V + \Omega_{T}(a)\times V$, and its field strength is $F_{T}(a\wedge b) = a\cdot\nabla\Omega_{T}(b) - b\cdot\nabla\Omega_{T}(a) + \Omega_{T}(a)\times\Omega_{T}(b)$. On the Einstein region $\Sch(A) = \tfrac{\Lambda}{6}A$, and the frozen point is $I = \nInf + \tfrac{\Lambda}{6}\nO$, $Q(I) = \Lambda/3$ (Fact \claimref{dark-energy.imported-frozen-point-geometry}); write $A = \hb^{-1}(a)$ for the physical frame vector.

(i) The commutator with a constant vector gives $D_{a}I = \Omega_{T}(a)\cdot I$; here $\Omega(a)\cdot I = 0$ (a spacetime bivector against a vector of the added plane), $(\nInf\wedge A)\cdot I = -\tfrac{\Lambda}{6}A$ and $(\nO\wedge\tfrac{\Lambda}{6}A)\cdot I = -\tfrac{\Lambda}{6}A$, so $D_{a}I = \tfrac{\Lambda}{6}A - \tfrac{\Lambda}{6}A = 0$: $I$ is parallel. Then $[D_{a}, D_{b}]I = F_{T}(a\wedge b)\cdot I = 0$, so at each point $F_{T}(a\wedge b)$ lies in the stabiliser algebra of $I$. Since $Q(I) \neq 0$ the vectors split as $\R I \oplus I^{\perp}$, and a bivector $B$ with $B\cdot I = 0$ is a bivector of $I^{\perp}$ (write $B = I\wedge u + B'$ with $u, B'$ in $I^{\perp}$; then $B\cdot I = -Q(I)\,u$). Transport by the rotor of $\Omega_{T}$ fixes $I$, hence preserves $I^{\perp}$ and the induced product, whose signature is $(1,4)$ or $(2,3)$ according to the sign of $\Lambda$ by Fact \claimref{dark-energy.imported-frozen-point-geometry}; so the holonomy lies in $SO(1,4)$ or $SO(2,3)$ and $D$ restricts to a metric connection on $I^{\perp}$.

(ii) By Lemma \claimref{dark-energy.stelle-west-vector-frozen-point}(i), $V = X - \tfrac{3}{\Lambda}I = \tfrac12\nO - \tfrac{3}{\Lambda}\nInf$ has $V\cdot I = 0$ and $Q(V) = -3/\Lambda$, and $I\wedge V = \nInf\wedge\nO = e_{+}\wedge e_{-}$ for either sign of $\Lambda$. The subgroup fixing $I$ and $V$ acts on $\{I, V\}^{\perp}$, which is the spacetime $\Cl(1,3)$ itself (both vectors lie in the added plane), of signature $(2,4) - (1,0) - (0,1) = (1,3)$ for $\Lambda > 0$ and $(2,4) - (0,1) - (1,0) = (1,3)$ for $\Lambda < 0$, and is its orthogonal group. The covariant derivative of $V$ is $D_{a}V = D_{a}X = -(\nInf\wedge A)\cdot\nO = A = \hb^{-1}(a)$, the soldering form (Lemma \claimref{dark-energy.stelle-west-vector-frozen-point}(ii)). The reduced connection is $\Omega_{T}$ itself: with $V/\ell^{2} = \tfrac{\Lambda}{3}V = \tfrac{\Lambda}{6}\nO - \nInf$,
\[
\Omega_{T}(a) = \Omega(a) - \bigl(\nInf - \tfrac{\Lambda}{6}\nO\bigr)\wedge A = \Omega(a) - \hb^{-1}(a)\wedge V/\ell^{2},
\]
the gauge field $D + e\wedge V/Q(V)$ of Lemma \claimref{dark-energy.stelle-west-vector-frozen-point}(iii), a bivector of $I^{\perp}$. Its field strength has Lorentz block $R(a\wedge b) + (A\wedge V)\times(B\wedge V)/\ell^{4}$ plus the translational term $(D_{a}B - D_{b}A)\wedge V/\ell^{2}$, which is the torsion and vanishes; and $(A\wedge V)\times(B\wedge V) = -Q(V)\,A\wedge B = \ell^{2}A\wedge B$ for $A, B$ orthogonal to $V$, so the Lorentz block is $R + \ell^{-2}e\wedge e$ (Lemma \claimref{dark-energy.stelle-west-vector-frozen-point}(iv)), in operator form $-(\mathcal{R}(B) - \ell^{-2}B)$ in the store's sign, zero on de Sitter. On an Einstein field it is $-\mathcal{W}$, a spacetime bivector, so it commutes with $I$ and $V$: the curvature takes values in the Lorentz algebra of $\{I, V\}^{\perp}$. This is the datum of Fact \claimref{dark-energy.imported-macdowell-mansouri-theorem-stelle-west} term by term.

(iii) The fibre volume form of $\mathcal{T}$ is $\epsilon(a, b, X, Y) = \langle a\wedge b\wedge X\wedge Y\,I_{6}^{-1}\rangle$ for vectors $a, b$ and bivectors $X, Y$, with $I_{6} = I_{4}E$, $E = e_{+}e_{-}$. With $I\wedge V = E$ and $X, Y$ spacetime bivectors, which commute with $E$,
\[
\epsilon(I, V, X, Y) = \langle E\,(X\wedge Y)\,E^{-1}I_{4}^{-1}\rangle = \langle (X\wedge Y)\,I_{4}^{-1}\rangle ,
\]
which is the Stelle-West contraction $\epsilon_{abcd}X^{ab}Y^{cd}$ of Fact \claimref{dark-energy.imported-macdowell-mansouri-theorem-stelle-west} (the six-dimensional form contracted with the unit normals $I/|I|$ and $V/\ell$, the normalising factors multiplying to $|Q(I)Q(V)|^{1/2} = 1$ on both branches, since $Q(I)Q(V) = (\Lambda/3)(-3/\Lambda) = -1$). So the Stelle-West prefactor $\ell^{2} = 3/\Lambda = 1/Q(I)$ is all that remains, up to the fixed constant of the functional: $S \sim Q(I)^{-1}\int\epsilon(I, V, F, F)$. The second form follows from $V = X - (3/\Lambda)I$ and the antisymmetry of $\epsilon$ in its first two slots, $\epsilon(I, V, \cdot, \cdot) = \epsilon(I, X, \cdot, \cdot)$. Slot count: six slots, two filled by $I$ and $V$, four by the two bivectors; the integrand is a scalar, which is why two vectors are needed and one is not enough. The field space: an ambient connection that does not preserve $I$ has components $D I \neq 0$, which enter through $[I\wedge u, I\wedge v] = -Q(I)\,u\wedge v$; the ambient rotors $R$ act on $(\Omega_{T}, I, V)$ together, and $R(B\times I)\tilde R = (RB\tilde R)\times(RI\tilde R)$ for a unit rotor, so the condition $D I = 0$ is carried to itself.

(1) is inspection: every ingredient of $\epsilon(I, V, F, F)$, and of its value $\epsilon(I, X, F, F)$ at the canonical solution, is a field of the ambient model or built from one, $F$ is polynomial in the gauge field, and no inverse of $\hb$ occurs; the locators for the continuation of the bundle and the connection across the wall are those of Fact \claimref{dark-energy.imported-walls-sphere-theorem}.

(2) The pole is the rational-function statement $1/Q(I) = 3/\Lambda$; polynomiality of the integrand in $\Lambda$ holds because $I = \nInf + \tfrac{\Lambda}{6}\nO$ and $\Omega_{T}$ are polynomial in $\Lambda$. The expansion is bilinearity: with $F = R + \ell^{-2}e\wedge e$, $\ell^{2}\epsilon(F, F) = \ell^{2}\epsilon(R, R) + 2\epsilon(R, e\wedge e) + \ell^{-2}\epsilon(e\wedge e, e\wedge e)$, which is the expansion of Fact \claimref{dark-energy.imported-macdowell-mansouri-theorem-stelle-west} with $\ell^{2} = 3/\Lambda$. The interior variation of the first term vanishes: $\delta R(a\wedge b) = D_{a}\delta\Omega(b) - D_{b}\delta\Omega(a)$, so $\delta\int\epsilon(R, R) = 2\int\epsilon(D\wedge\delta\Omega, R)$, which is a boundary term minus $2\int\epsilon(\delta\Omega, \dot D\wedge\dot R)$, and the second Bianchi identity $\dot D\wedge\dot R = 0$ (dictionary row 5) removes the latter. The remaining two terms are polynomial in $\Lambda$, and at $\Lambda = 0$ the cosmological term vanishes, leaving the Einstein-Cartan functional $2\int\epsilon(R, e\wedge e)$, whose only ingredients are the frame $\hb^{-1}$, the field strength and the pseudoscalar: no inverse of $\hb^{-1}$ is taken, so it is defined at degenerate frames.

\emph{The checks.} The exact checks compute all of this in $\Cl(2,4)$ on gauge fields: the stabiliser of $I$ (dimension $10$, $I^{\perp}$ of signature $(1,4)$ or $(2,3)$) and the tractor field strength with its derivative terms (zero on de Sitter, the Weyl curvature on the Kasner vacuum, annihilating $I$ and $V$); the Stelle-West datum on de Sitter; and, through the shared engine, on Kottler fields of both signs of $\Lambda$ in the Painlev\'e-Gullstrand chart, $\Sch = \tfrac{\Lambda}{6}\mathrm{id}$ from the field strength, $D I = 0$, $D_{a}V = \hb^{-1}(a)$, $\Omega_{T}(a) = \Omega(a) - \hb^{-1}(a)\wedge V/\ell^{2}$, $I\wedge V = e_{+}\wedge e_{-}$ and $Q(I)Q(V) = -1$, and the tractor field strength, computed with derivatives, equal to $-\mathcal{W}$ and commuting with $I$ and $V$; and, for (iii) and (2), $\epsilon(I, V, F, F) = \epsilon(I, X, F, F) = \epsilon(\nInf, \nO, F, F)$, the prefactor $3/\Lambda$, the expansion at four values of $\Lambda$, $[I\wedge u, I\wedge v] = -Q(I)\,u\wedge v$, and the ambient invariance under a sample rotor.

\emph{In components} (a labelled reading). In the splitting of the derived metric the computation of (ii) is the index form of Lemma \claimref{dark-energy.stelle-west-vector-frozen-point}: the bracket of the translational part with itself is $\tfrac12\cdot 2\cdot(\Lambda/3)\,e\wedge e$ in the Lorentz block, that is $\ell^{-2}e\wedge e$, and the mixed terms $[D, e\wedge V]$ carry the torsion $De$ into the translational block. In forms, $\epsilon_{abcd}F^{ab}F^{cd} = \epsilon(I, V, F, F)/(|I|\,\ell)$, and the variation of (2) reads $\delta\int\epsilon_{abcd}R^{ab}\wedge R^{cd} = 2\int d\bigl(\epsilon_{abcd}\,\delta\omega^{ab}\wedge R^{cd}\bigr) - 2\int\epsilon_{abcd}\,\delta\omega^{ab}\wedge DR^{cd}$ with $DR = 0$.
